% !TeX root = ../main.tex
\section{Task Definition and Evaluation}
\label{sec:task}

\subsection{Outputs}

A tone estimate can take three forms. (i) An \emph{interpretable chain}: an
ordered list of effect types with parameter values, which a player can load
and adjust. (ii) A \emph{neural model} of the processing, such as an
amplifier capture, which reproduces the sound but exposes few or no
controls. (iii) \emph{Processed audio} only, as produced by tone-transfer
models that re-render a recording in another tone. Only (i) and (ii) are
usable in a real-time effects application, and (i) is the one players can
edit. Sub-tasks studied separately in the literature are effect detection
(which effects are present), chain order, parameter regression for a known
device, device identification (which amplifier or pedal), and estimation of
the linear part alone (an equalisation curve or an impulse response).

\subsection{Settings by available signals}

Methods differ most in what they require. \emph{Paired}: the dry signal
(direct input, DI) and the processed signal of the same performance, aligned
sample by sample; this is how amplifiers are captured. \emph{Unpaired}: a
processed recording plus dry guitar from another performance, possibly on
another instrument. \emph{Wet only}: the processed recording alone. In the
commercial case the target is wet only and, in addition, separated from a mix;
the player's own guitar can supply unpaired dry audio.

\subsection{Metrics}

\emph{Parameter metrics} (classification accuracy, F1 over effect sets,
exact-chain accuracy, mean absolute error of knob values on a unit scale) need
ground truth and are meaningful only when the target was produced with the
same chain the estimator controls. Peladeau and Peeters found that training on
an audio loss gives better reconstructions than training on parameters even
when the parameter error is larger~\cite{peladeau2024blind}, so parameter error
should not be the main criterion.

\emph{Aligned audio metrics} compare the re-rendered signal with the target
sample by sample or frame by frame: error-to-signal ratio (ESR) with
pre-emphasis~\cite{wright2020perceptual}, multi-resolution STFT distance
(MR-STFT)~\cite{yamamoto2020parallel,steinmetz2020auraloss}, mel or MFCC
distances, SI-SDR. They need the same performance on both sides, so without a
DI they apply only to synthetic tests or to a recovered dry signal.

\emph{Embedding and distribution metrics} compare signals without alignment:
distances in an effect embedding such as AFx-Rep~\cite{steinmetz2024stito} or
Fx-Encoder++~\cite{yeh2025fxencoder}, CLAP~\cite{wu2023laionclap}, and Fréchet
Audio Distance (FAD)~\cite{kilgour2019fad}. FAD is biased by sample size and
depends on the embedding and the reference set; Gui et al.\ propose an
extrapolated FAD$_\infty$~\cite{gui2024adapting}, and the kernel variant KAD
converges with fewer samples~\cite{chung2025kad}. A single song gives few
clips, so pairwise embedding distances suit per-song evaluation better.

\emph{Listening tests} use MUSHRA (ITU-R BS.1534-3)~\cite{itu2015bs1534}, often
run with webMUSHRA~\cite{schoeffler2018webmushra}, or mean opinion scores split
into tone similarity and audio quality~\cite{yeh2026egvae}.

\subsection{Which metrics agree with listeners}

Two studies measured it for guitar devices (Table~\ref{tab:metrics}). In the
ToneTwist evaluation of 40 nonlinear devices (27 screened listeners),
MR-STFT had the strongest Spearman correlation with MUSHRA ratings among
signal metrics and ESR one of the weakest; among embedding-based FAD, VGGish
and AFx-Rep did best overall and CLAP best on distortion devices, while no
metric was reliable for overdrive or compression~\cite{comunita2025differentiable}.
Wright, Dal Santo and Schlecht collected ratings for 590 emulation--reference
pairs from published amplifier and distortion models and found MFCC and deep
embedding distances predictive, and MAE and ESR poor~\cite{wright2026predicting}.
In unpaired amplifier modelling, ESR and listener ratings disagreed on a
heavy-distortion tone~\cite{wright2023adversarial}. No metric has been
validated for delay, reverberation or modulation on guitar. Perceptual
ceilings are also informative: listeners barely told Kemper profiles from the
original amplifiers ($d' = 0.34$, $N = 177$)~\cite{duvel2020confusingly}.

\begin{table}[t]
\centering
\caption{Spearman correlation of objective metrics with MUSHRA ratings on 40
nonlinear guitar and bass devices~\cite{comunita2025differentiable}
(negative: lower distance, higher rating).}
\label{tab:metrics}
\footnotesize
\begin{tabular}{@{}lrlr@{}}
\toprule
Signal metric & $\rho$ & FAD embedding & $\rho$ \\
\midrule
MR-STFT & $-0.427$ & VGGish & $-0.430$ \\
MAPE & $-0.339$ & AFx-Rep & $-0.417$ \\
L1 & $-0.332$ & CLAP & $-0.406$ \\
MSE & $-0.314$ & PANN & $-0.303$ \\
ESR & $-0.311$ & & \\
\midrule
\multicolumn{4}{@{}l@{}}{High-distortion devices: MR-STFT $-0.536$,} \\
\multicolumn{4}{@{}l@{}}{FAD with CLAP $-0.613$.} \\
\bottomrule
\end{tabular}
\end{table}

\subsection{Comparability}

Most results use clean, isolated, plugin-rendered guitar, so they measure an
easier problem than the commercial one. Many test sets share the dry
recordings or the plugins of the training set. And because settings can be
non-unique, two estimates with very different parameters may sound alike,
which parameter metrics penalise and listeners do not.
